% einstein_photoelectric_ML_short.tex
% Short newsletter version of einstein_photoelectric_ML.tex
% Thomas J. Sargent, 2026

\documentclass[11pt]{article}

\usepackage{amsmath,amssymb}
\usepackage{booktabs}
\usepackage[margin=1.15in]{geometry}
\usepackage{microtype}
\usepackage{hyperref}
\usepackage{natbib}
\usepackage{parskip}

\hypersetup{colorlinks=true, linkcolor=blue, citecolor=blue, urlcolor=blue}

\title{Five Data Points and a Constant of Nature\\[2pt]
       \large Einstein's Photoelectric Equation as a Regression, 1905--1916}
\author{Thomas J.\ Sargent\thanks{%
  This paper is one of five that read episodes in the
  history of physics and astronomy as exercises in statistical inference and
  machine learning.  \citet{sargent2025sources} sets out the theme: what is
  now called artificial intelligence names practices that scientists have used
  for centuries.  The other four treat Kepler and Newton, Pauli's exclusion principle, the
  exoplanet discovery of Mayor and Queloz, and the accelerating universe of
  Riess, Schmidt and Perlmutter.}}
\date{March 2026}

\begin{document}
\maketitle

\begin{abstract}
\noindent
In 1905 Einstein wrote down a linear equation relating a measurable voltage to
the frequency of light, with a slope fixed in advance by Planck's blackbody
constant.  In 1916 Millikan measured the slope with five observations and got
it right to under one percent.  I run the regression, price its standard
error, and ask why so few observations sufficed.  The answer is that the
physics had already fixed the functional form, the number of free parameters,
and the units of the slope.  A learner starting from the data alone would have
needed thousands of observations and would still not have known what it had
found.  A shorter lesson also emerges: the regression settled the equation and
not the physics behind it, and the physics waited another sixty years for a
different measurement.  A longer version of this note carries the derivations
and the full data.
\end{abstract}

%======================================================================
\section{The Setup}

In December 1900 Max Planck introduced the energy quantum $h\nu$ to fit the
blackbody spectrum.  He called the step ``an act of desperation.''  Five years
later Einstein borrowed the quantum and applied it to a second phenomenon.

Shine light on a metal in a vacuum and electrons come off.  Apply a retarding
voltage and you can stop them.  The voltage that just stops the fastest
electron is the \emph{stopping potential} $V_0$.  Einstein proposed that light
arrives in discrete quanta of energy $h\nu$, that one quantum is absorbed by
one electron, and that the electron escapes if $h\nu$ exceeds the energy $W$
binding it to the metal.  Energy conservation then gives
\begin{equation}\label{eq:einstein}
  V_0 = \frac{h}{e}\,\nu - \frac{W}{e},
\end{equation}
where $e$ is the electron's charge.  A line.  Slope $h/e$, intercept $-W/e$.

Three features of \eqref{eq:einstein} would interest anyone who estimates
equations for a living.

The slope is a universal constant.  It does not depend on which metal you use.
Every metal must lie on a line of the same gradient, differing only in
intercept.  Economists will recognise a cross-equation restriction: one deep
parameter is required to appear identically in the reduced form of every
experiment.

The slope was known before the experiment.  Planck had already fitted $h$ to
the blackbody spectrum.  Einstein's equation therefore made a prediction with
no free parameter to adjust, which is a rarer thing than a good fit.

The intercept is contaminated.  Between the metal emitting the electrons and
the electrode collecting them sits a contact electromotive force $v_c$, a
voltage difference of purely instrumental origin.  It could reach 3 volts,
comparable to the stopping potentials themselves, and it drifted with surface
age and adsorbed gas.  It enters the observed reading additively:
\begin{equation}\label{eq:obs}
  V_0^{\rm obs} = \underbrace{-W/e - v_c}_{\text{intercept}}
    + \frac{h}{e}\,\nu + \varepsilon .
\end{equation}
A nuisance parameter that shifts the level and leaves the slope alone.  The
slope is identified from variation in $\nu$; the intercept is not identified
without an independent measurement of $v_c$.

One more thread runs out of the 1905 paper, and it took twenty years to
surface.  Einstein's route to the light quantum was thermodynamic.  He showed that in one regime the entropy of
radiation depends on volume as the entropy of an ideal gas of $n$ independent
particles does.  He inferred that radiation behaves, for those purposes, as a
collection of independent localised quanta.  The step
from there to treating the classical wave intensity as a probability density
for finding a quantum came later, in the 1920s, when Einstein called the wave
a \emph{Gespensterfeld}, a ghost field guiding probabilities without carrying
energy.  Max Born generalised the idea from photons to the Schr\"odinger wave
function in 1926 \citep{Born1926} and made it the interpretive core of quantum
mechanics.  Einstein won the 1921 Nobel Prize for the photoelectric equation.
Born won the 1954 Prize for the probabilistic reading that the equation's
underlying picture had suggested.  Einstein spent the rest of his life
opposing it.

\subsection*{The harder case, in the note that follows}

The note that follows applies this anatomy to the discovery of cosmic
acceleration by Riess, Schmidt and Perlmutter in 1998.  General relativity
fixes the functional form there, in place of the quantum hypothesis.  A
combination of the Hubble constant and the supernovae's absolute magnitude
plays the part of Millikan's contact voltage, shifting the level of the
diagram and leaving its shape alone.  The fit again settles a functional form
without settling the physics behind it.

The photoelectric case is the easy one, and worth taking first for that
reason.  Millikan wanted a single number, and removing the nuisance parameter
left that number identified.  The Friedmann model has two shape parameters.  Removing the nuisance
parameter there leaves them entangled with each other, and the second note
takes up what becomes of identification when it does.

%======================================================================
\section{Why It Took Eleven Years}

Einstein's equation was published in 1905 and confirmed in 1916.  The delay
was not a failure of statistics.  Three groups tried in between and got
answers scattered over a factor of two.

\citet{Hughes1913} used three ultraviolet lines spanning 69 nanometres and
about a dozen metals, and reported values of $h$ ranging from
$3.55$ to $5.85 \times 10^{-27}$ erg$\cdot$s.  \citet{RichardsonCompton1912}
measured eight metals and did no better.  Every pre-1916 estimate was low by
between 12 and 47 percent.

A discrepancy of that size is not sampling error.  Three systematic
mechanisms produced it.  Stray short-wavelength light raised the apparent
$V_0$ at long-wavelength lines and flattened the slope.  Contact potentials
drifted between sessions, contaminating observations one at a time instead of
shifting a common intercept.  And the frequency range was too narrow, which
inflated the standard error and left the estimate hostage to a single bad
line.

Only the third of these is a statistical problem.  Estimation theory cannot
repair the other two after the fact.

Robert Millikan spent a decade building an apparatus that removed all three.
An electromagnetically driven knife shaved a fresh sodium surface inside the
vacuum.  Filters passed only mercury-arc lines with no short-wavelength
companions.  A Kelvin double-balance monitored $v_c$ continuously.  And the
spectral lines spanned 293 nanometres rather than 69.

%======================================================================
\section{The Regression}

Millikan's sodium data, corrected for the contact EMF he measured
independently, are in Table~\ref{tab:data}.

\begin{table}[ht]
\centering
\caption{Millikan's sodium stopping potentials, corrected for a contact EMF
of $2.51$~V measured by the Kelvin method.  Five of the six values are
reconstructions read off the line Millikan fitted to his Figure~5 and rounded
to three decimals; only the 2535~\AA\ entry is a directly printed figure.}
\label{tab:data}
\medskip
\begin{tabular}{rrr}
\toprule
$\lambda$ (\AA) & $\nu$ ($10^{14}$ Hz) & $V_0$ (V) \\
\midrule
5461 &  5.490 & 0.454 \\
4339 &  6.908 & 1.039 \\
4047 &  7.409 & 1.246 \\
3650 &  8.213 & 1.578 \\
3126 &  9.591 & 2.147 \\
2535 & 11.827 & 3.030 \\
\bottomrule
\end{tabular}
\end{table}

Ordinary least squares on the first five points gives
\begin{equation}\label{eq:fit}
  \hat\beta = \frac{S_{\nu V}}{S_{\nu\nu}} = \frac{3.830}{9.277}
    = 4.128\times10^{-15}~\text{V}\cdot\text{s},
  \qquad \hat\alpha = -1.813~\text{V},
\end{equation}
where $S_{\nu\nu} = \sum_i (\nu_i - \bar\nu)^2$ measures how widely the
frequencies are spread and $S_{\nu V} = \sum_i (\nu_i - \bar\nu)(V_{0i} -
\bar V_0)$ their covariation with the voltages.  $S_{\nu\nu}$ carries the whole
design argument of Section~\ref{sec:design}.

The residuals from this fit are all smaller than a quarter of a millivolt,
forty times below Millikan's stated reading precision of $\pm10$~mV.  That is
not evidence of a very precise experiment.  The reconstruction noted in
Table~\ref{tab:data} produced it.  Five of these numbers were read off a
fitted line and rounded to three decimals, so the residuals measure
rounding.  Taken at face value they give a
relative standard error of $0.0015\%$, three hundred times finer than Millikan
himself claimed.

The defensible procedure takes $\sigma$ from the instrument rather than from
the residuals.  Millikan's balance points were reproducible to about
$10$~mV, and
\begin{equation}\label{eq:se}
  \mathrm{SE}(\hat\beta) = \frac{\sigma}{\sqrt{S_{\nu\nu}}}
  = \frac{0.010}{\sqrt{9.277}\times10^{14}} = 3.3\times10^{-17}~
  \text{V}\cdot\text{s},
\end{equation}
a relative uncertainty of $0.8\%$ for the five-line design and $0.5\%$ for the
six-line design.

Millikan also determined the slope nine separate times, pairing one spectral
line with 5461~\AA\ in nine sessions with a freshly shaved surface each time.
Those nine values are genuine replications.  Their mean is
$4.131\times10^{-15}$, their standard deviation $0.087$, and the standard
error of the mean $0.029$, or $0.71\%$.

Two routes, one from replication and one from the information matrix, agree to
within a tenth of a percent.  That agreement, and not the spurious $0.0015\%$,
is what the uncertainty on $h$ rests on.

\subsection*{The cross-metal restriction}

Because $h/e$ is universal, every metal must give the same slope.  Millikan
also ran lithium.  Fitting five lithium lines gives
$\hat\beta_{\rm Li} = 4.127\times10^{-15}$, differing from the sodium slope by
$0.04\%$.  Pooling the two metals with a common slope and separate intercepts
gives $4.127\times10^{-15}$ with a standard error of $0.49\%$, since pooling
enlarges $S_{\nu\nu}$ in the same way that adding the sixth sodium line does.
The Wald statistic for equality of the two slopes is $0.04$.

The restriction is not rejected, and it would be a mistake to call that a
confirmation.  The test has little power.  With a standard error of $0.8\%$ on
each slope it would fail to reject differences of up to about $2\%$, so it
excludes only gross violations.  More seriously, four of the five lithium
entries in Millikan's table are the sodium entries shifted by a constant
$0.540$~V, and a constant shift leaves the slope unchanged by construction.
The near-zero test statistic reflects the arithmetic as much as the physics.
Millikan established that the two metals' curves were parallel at his stated
precision, which is how he put it.

%======================================================================
\section{Design Beats Sample Size}
\label{sec:design}

Under Gaussian errors, maximising the likelihood and minimising the sum of
squared residuals are the same operation, so the least-squares estimates in
\eqref{eq:fit} are also maximum likelihood estimates.  The Fisher information
for $(\alpha,\beta)$ is $\sigma^{-2}X^\top X$, and inverting it gives the
Cram\'er--Rao bound on the variance of any unbiased estimator of the slope,
\begin{equation}\label{eq:crb}
  \operatorname{Var}(\tilde\beta) \;\geq\; \frac{\sigma^2}{S_{\nu\nu}},
  \qquad S_{\nu\nu} = \sum_i (\nu_i - \bar\nu)^2 ,
\end{equation}
which least squares attains exactly.  Nothing can do better.

Equation \eqref{eq:crb} is a statement about the \emph{design} and not about
the readings.  $S_{\nu\nu}$ depends only on which spectral lines the
experimenter chose.  Spreading the frequencies apart shrinks the bound.
Collecting more readings at the same frequencies does much less.

\begin{table}[ht]
\centering
\caption{Effect of the frequency range on the standard error of $h/e$, holding
the instrumental $\sigma = 10$~mV fixed.  The narrow window is the closest
analogue Millikan's line set permits to the bands available to Hughes and to
Richardson \& Compton.}
\label{tab:design}
\medskip
\begin{tabular}{lccrr}
\toprule
Design & $n$ & Span (nm) & $S_{\nu\nu}$ ($10^{28}$ Hz$^2$) & Rel.\ SE \\
\midrule
Narrow (4047--3126 \AA)        & 3 & 92  & 2.44  & 1.55\% \\
Millikan 5-point               & 5 & 234 & 9.28  & 0.79\% \\
Millikan 6-point               & 6 & 293 & 24.72 & 0.49\% \\
\bottomrule
\end{tabular}
\end{table}

Going from three lines to six cuts the standard error by a factor of three,
and almost all of the gain comes from the wider span rather than from the
extra observations.  The point generalises past photoelectrics.  Where a parameter
is identified by variation in a regressor, the design of that variation
dominates the sample size.

The sixth line illustrates the cost of the same strategy.  Adding the
2535~\AA\ point, the only directly quoted number in the table, moves the slope
to $4.070\times10^{-15}$, $1.4\%$ down.  It sits $0.04$~V below the five-point
line, and its extreme frequency gives it the highest leverage.  Millikan
attributed the deficit to scattered ultraviolet light and dropped the line
from his primary determination.  The observation that contributes most to
$S_{\nu\nu}$ is also the one most exposed to systematic error.

%======================================================================
\section{The Prediction}

Converting the slope to Planck's constant needs the electron charge.  With the
value Millikan had measured himself, $e = 1.592\times10^{-19}$~C, the
five-point slope gives $h = 6.573\times10^{-34}$~J$\cdot$s, low by $0.81\%$
against the modern defined value.  With the modern $e$ it gives
$6.616\times10^{-34}$, low by $0.17\%$.  Millikan's error came from his
oil-drop charge, which was too small by $0.64\%$, and not from the
photoelectric slope his 1916 apparatus was built to measure.

The sharper test does not involve the modern value at all.  Planck had fitted
$h = 6.55\times10^{-27}$~erg$\cdot$s to the blackbody spectrum in 1901.  With
Millikan's own $e$, that predicts a slope
\begin{equation}\label{eq:pred}
  \left(\frac{h}{e}\right)_{\rm predicted}
  = \frac{6.55\times10^{-34}}{1.592\times10^{-19}}
  = 4.114\times10^{-15}~\text{V}\cdot\text{s}.
\end{equation}
Millikan measured $4.128\times10^{-15}$.  The two agree to $0.35\%$, inside
the $0.8\%$ standard error.

No photoelectric measurement entered the prediction.  Planck fitted $h$ to a
different phenomenon fifteen years earlier.  Einstein derived the slope in
1905.  Millikan measured it in 1916.  The regression had one free parameter,
the intercept, and it was not the one being tested.

Both 1916-era constants were too small, and the errors partly cancel in the
ratio.  Planck's $h$ was low by $1.15\%$ and Millikan's $e$ by $0.64\%$, so
the predicted slope was low by $0.52\%$ and the measured slope by $0.17\%$.
Against modern constants the comparison is $4.128$ against $4.136$.

%======================================================================
\section{The Equation Is Still in Service}

Equation \eqref{eq:einstein} did not become a historical curiosity.
Laboratories with no interest in Planck's constant read it backwards every
day.

In X-ray photoelectron spectroscopy the incident photon energy is known and
the measured electron energy identifies the binding energy of the core
electron the photon knocked loose:
\begin{equation}\label{eq:xps}
  E_B = h\nu - eV_0 - \phi_{\rm spec},
\end{equation}
where $\phi_{\rm spec}$ is an instrument calibration constant.  That is
\eqref{eq:einstein} rearranged, with a different unknown.  Because every
element has its own set of core-level binding energies, the technique reads
the elemental composition of a surface directly.  Kai Siegbahn received the
1981 Nobel Prize for it \citep{Siegbahn1967}, and it is now standard in
semiconductor fabrication, catalysis and materials science.  Ultraviolet
photoelectron spectroscopy and Kelvin probes apply the same equation to
measure work functions of semiconductor films and photovoltaic layers to a few
meV.

Einstein's equation survives in quantum electrodynamics unchanged, as the
energy balance for the fastest photoelectron: the electron that starts at the
Fermi level and reaches the detector without inelastic loss.  The balance is exact.  The measured
spectrum is not.  A real metal broadens the emission
edge thermally, the work function belongs to a crystal face and its adsorbate
layer rather than to an element, and final-state relaxation shifts core levels
by amounts that matter at spectroscopic precision.  Millikan's residual
scatter came from the surface physics.  The energy balance was never the
problem.

%======================================================================
\section{The Man Who Confirmed It Did Not Believe It}

Millikan's 1916 paper confirmed Einstein's equation to better than one
percent.  In the same paper he wrote:
\begin{quote}
  ``Despite then the apparently complete success of the Einstein equation, the
  physical theory of which it was designed to be the symbolic expression is
  found so untenable that Einstein himself, I believe, no longer holds to
  it.'' \citep[p.\ 388]{Millikan1916}
\end{quote}

He had built the apparatus hoping to \emph{refute} the equation.  The
light-quantum hypothesis struck him as physically impossible, since a stream of
localised lumps could not diffract or interfere, and diffraction and
interference were not in doubt.  He got the opposite of what he expected and
published it anyway, with the interpretation disowned.

The stance is more defensible than it sounds, and it will be familiar to
anyone who has estimated a reduced form.  Millikan confirmed a predictive
relationship without endorsing the structural model that generated it.  He was
right to separate the two, and Section~\ref{sec:limits} shows that he was
right on the merits as well: the regression really does not establish the
mechanism.  Where
he erred was in thinking the equation's success left the hypothesis
untouched.

Compton scattering \citep{Compton1923} supplied independent evidence for the
particulate character of light in 1923.  Millikan came to describe his own
result differently:
\begin{quote}
  ``This work resulted, contrary to my own expectation, in the first direct
  experimental proof \ldots\ of the complete validity of the Einstein
  equation.''
\end{quote}
\noindent The sentence belongs to his later retrospective writing rather than
to the Nobel lecture he delivered in 1924.  \citet{Holton1999} and
\citet{Pais1982} trace how his views moved.

%======================================================================
\section{Einstein as Machine Learner}

Recast the episode in the vocabulary of statistical learning.  The
translation is not a conceit.  \citet{sargent2025sources} argues that the
techniques now gathered under artificial intelligence have long histories
inside the sciences that used them, and that reading old episodes this way
shows what the methods can and cannot do.

The \emph{hypothesis class} is the set of functions the estimator is allowed to
consider.  Einstein's is
\begin{equation}\label{eq:class}
  \mathcal{H} = \{ \nu \mapsto \alpha + \beta\nu \},
\end{equation}
with $\beta$ shared across metals.  Two free parameters per metal, one shared
slope, pseudo-dimension two.  The restriction was not a modelling convenience.
It followed deductively from $E = h\nu$ and energy conservation.  The quantum
hypothesis \emph{is} the inductive bias.

The \emph{feature selection} was Lenard's.  In 1902 he established that the
stopping potential does not depend on light intensity and does depend on
frequency.  Classical wave theory predicted the opposite.  Lenard's experiment
removed intensity from the feature set before any line was fitted.

The \emph{training} is least squares, and for this class it has a closed form.
No gradient descent, no learning rate, no early stopping.

Now compare the alternative.  A three-layer network with sixteen ReLU units
per layer has
\[
  p = (1\!\cdot\!16+16)+(16\!\cdot\!16+16)+(16\!\cdot\!16+16)+(16\!\cdot\!1+1)
    = 593
\]
free parameters and a pseudo-dimension near $5{,}500$.  The pseudo-dimension
is the capacity measure appropriate to real-valued function classes; for a
linear model it equals the number of free parameters, which is why Einstein's
class scores two.  The practitioner's rule of thumb, $n \gtrsim 10p$, calls for
about six thousand training examples.  Millikan had five.

\begin{table}[ht]
\centering
\caption{What was available when, against what each approach needs.  The
network's requirement uses $n \gtrsim 10p$; the formal uniform-convergence
bound is four orders of magnitude larger and vacuous here.}
\label{tab:complexity}
\medskip
\begin{tabular}{lrll}
\toprule
Epoch & $n$ & Einstein's model & Neural network \\
\midrule
Lenard 1902   & $\approx 0$ & Sufficient$^\dagger$ & Cannot train \\
Hughes 1912   & $\sim 30$   & Sufficient & Short by $\times$200 \\
Millikan 1916 & 5           & Sufficient & Short by $\times$1200 \\
Modern era    & $\sim 5000$ & Sufficient & Short by $\times$1.2 \\
\bottomrule
\end{tabular}

\smallskip
\begin{minipage}{0.85\textwidth}
\footnotesize $^\dagger$ Theory fixes the functional form; two parameters
remain to be estimated from data.
\end{minipage}
\end{table}

The last row is the interesting one.  Given enough clean observations a
network would eventually trace out the line.  It would still not identify the
fitted gradient as the ratio $h/e$ of two independently measured constants of
nature.  It would still not impose the cross-metal restriction from first
principles, seeing universality as a coincidence to be discovered rather than
a constraint to be applied.  And it would not decompose the intercept into a
work function and an apparatus nuisance, because those two are not separately
identified without the external Kelvin measurement.

Identification failures of this kind are invisible to a fitting procedure.  The
regression runs, reports a number, and says nothing about what the number
means.

Table~\ref{tab:loop} sets the two vocabularies side by side.

\begin{table}[ht]
\centering
\caption{The photoelectric episode in two vocabularies.}
\label{tab:loop}
\medskip
\begin{tabular}{p{4.3cm}p{3.6cm}p{4.3cm}}
\toprule
Step & Scientific name & Machine-learning name \\
\midrule
Propose $E = h\nu$ & Theoretical hypothesis & Model architecture \\
Derive $V_0 = (h/e)\nu - W/e$ & Deduction & Hypothesis class \\
Use $\nu$, not intensity & Feature selection & Feature engineering \\
Measure $(\nu_i, V_{0i})$ & Experiment & Training data \\
Minimise $\sum \hat\varepsilon_i^2$ & Least squares & Squared-error loss \\
Slope matches Planck's $h/e$ & Prediction confirmed & Out-of-sample test \\
Na and Li share a slope & Universality & Parameter tying \\
\bottomrule
\end{tabular}
\end{table}

Five points sufficed because the physics had already supplied the functional
form, the parameter count, and the units of the slope.  Theory does not
substitute for data.  It compresses how much data you need.

\subsection*{Where the compression runs out}

The compression is not uniform across the parameters, and the intercept shows
where it fails.  Millikan's fitted intercept implies a sodium work function of
$-\hat\alpha = 1.81$~eV.  The modern value for clean bulk sodium is $2.36$~eV.

Part of the gap is estimation error.  The intercept extrapolates to $\nu = 0$,
far outside the observed range, and its standard error is about $0.08$~V at
$\sigma = 10$~mV, an order of magnitude worse in relative terms than the
slope's.  The rest is physics.  Millikan's freshly shaved surfaces sat in
residual gas, and adsorbed gas shifts a metal's work function by several
tenths of a volt.  The work function also varies by some $0.3$~eV across the
crystal faces of a single metal.  Sodium has no one work function for the
intercept to be compared against.

The two coefficients are identified on different terms.  The slope is a
universal constant, identified by variation within a run, and immune to any
disturbance common to all the lines.  The intercept is a surface-specific
quantity, identified by extrapolation, and contaminated by exactly the
conditions the apparatus could not control.  Einstein's equation is far better
identified in its slope than in its intercept.  The slope carries the
physics.

%======================================================================
\section{What the Regression Could Not Settle}
\label{sec:limits}

The fit is excellent and the prediction was confirmed.  Neither establishes
that light is made of photons.

\citet{Lamb1969} showed that equation \eqref{eq:einstein} follows from a
\emph{semiclassical} model in which the electromagnetic field stays a
classical wave and only the atom is quantised.  The equation is a statement of
energy conservation together with quantised \emph{atomic} energy levels.  It
does not require the radiation field itself to come in lumps.

So Millikan's regression selected a functional form and ruled out a rival
form.  It could not choose between two physical stories that imply the same
form.  A model-selection criterion scoring in-sample fit would have reported
overwhelming evidence and would have been answering a different question from
the one everyone thought was being asked.

Settling it took a different observable and sixty more years.  The decisive
evidence is photon \emph{antibunching}: photons from a single quantum emitter
arrive \emph{less} bunched than Poisson statistics allow.  No classical
stochastic field can produce sub-Poissonian counts, however it is modulated.
\citet{Kimble1977} observed the effect in the resonance fluorescence of
sodium atoms.

The experiment was designed to test field quantisation, and the observable it
measured, photon-number variance, follows from Glauber's quantum theory of
optical coherence.  Without that theory nobody would have thought to measure
photon-number correlations at all.

The photoelectric episode is usually told as a confirmation.  A division of
labour describes it better.

%======================================================================
\section{What This Says About AI in the Sciences}

Deep networks are the right tool for problems in which no parsimonious
parametric family is known.  Image recognition has no closed form.  Neither
does protein structure prediction, and AlphaFold's achievement was to learn a
map that six decades of physical modelling had failed to write down.  Neither
does language.  In each case the hypothesis class is genuinely unknown, the
input is high-dimensional, and data are abundant.  Learning the function is
the only available move.

The photoelectric problem sits at the opposite corner of the design space: one
input, one output, five observations, and a functional form already derived
from a physical principle.  Nothing about the success of learned models on the
first class of problem bears on the second.

Two claims often made for AI in the sciences are worth separating against this
example.

The first is that a sufficiently capable learner can rediscover physical laws
from data.  In a limited sense this is true and has been demonstrated.
Symbolic regression systems such as AI Feynman \citep{Udrescu2020} recover the
functional forms of textbook relations from tabulated data.  But they are
given the units and the physical constants as inputs, and it is the units that
let a fitted coefficient be matched to $h/e$ rather than reported as an
unexplained number.  That hardcoding is prior physical knowledge.  The
$(\nu, V_0)$ data do not contain it.

The second is that scaling data collection substitutes for theory.  The
Cram\'er--Rao bound in \eqref{eq:crb} says otherwise, and says it precisely.
The standard error falls with the \emph{spread} of the regressor, not with the
count of observations at fixed design.  Millikan could have taken ten thousand
readings inside the 92-nanometre window and done worse than five readings
across 293 nanometres.  Deciding where to look is a modelling decision, and
the model that tells you where to look is the thing being tested.

The sharper limit is the one in Section~\ref{sec:limits}.  Even a perfect fit
of a correctly
specified equation to unlimited data would not have settled whether light is
quantised, because a rival physical theory implies the same equation.  Getting
that question answered required somebody to prove that the two theories differ
in photon-number variance, and then somebody to measure photon-number
variance.  No quantity of $(\nu, V_0)$ pairs contains the answer.  A system
that optimises fit on the data it is given cannot discover that it is
measuring the wrong thing.

%======================================================================
\section{Summary}

Einstein's photoelectric equation is a two-parameter regression whose slope
was fixed in advance by a constant fitted to an unrelated phenomenon.  Five
observations measured it to under one percent.  The regression itself is
trivial; obtaining clean data took a decade of instrument building, and
widening the range of the regressor mattered more than adding observations.
The final estimates are
\[
  \hat\beta = 4.128\times10^{-15}~\text{V}\cdot\text{s} \pm 0.8\%,
  \qquad
  h = 6.62\times10^{-34}~\text{J}\cdot\text{s},
  \qquad
  W_{\rm Na}/e = 1.81~\text{eV}.
\]

Two results point in opposite directions.  A correctly specified model turned
a problem needing thousands of observations into one needing five.  The same
model, fitted to those five points as precisely as anyone could wish, left the
physical question open for another sixty years, until a theory arrived that
named a different thing to measure.

\bigskip
\begin{center}\rule{0.4\textwidth}{0.4pt}\end{center}
\smallskip
\noindent\footnotesize
Millikan's tabulated stopping potentials, apart from the 2535~\AA\ entry, are
reconstructions rather than independently recovered readings, and the standard
errors quoted here come from his instrumental precision and his replicated
sessions rather than from the residuals of the fitted line.  Planck's 1901
value for $h$ and the sourcing of the Millikan quotations come from secondary
sources and should be checked against the originals.

\bigskip
\bibliographystyle{plainnat}
\bibliography{einstein_photoelectric_ML}

\end{document}
